% !TeX root = ../papers/nc-p3.tex
\section{種数二の核曲線と余分な三点繊維}
\label{sec:genus-two-kernel-excess}
\begin{theorem}\label{thm:genus-two-kernel-excess}
被約四辺有限正則窓の整な反対八次で、核像が曲線のとき、その余因子スキームの
重複度込みの曲線次数は十八以下である。階数一曲線が存在しても成立する。
\end{theorem}
\begin{proof}[種数二曲線自体への降下]
定理\ref{thm:sectional-cofactor-length}の等号 $d=19$ を仮定する。
一般平面の自己準同型曲線 $G$ は核空間内に五次曲線として閉埋込みされる。
その像は核像 $K$ 全体であるから $K=G$ で、特に整Gorenstein、算術種数二である。
$P=\mathcal O_K(1)$ とする。定理\ref{thm:kernel-curve-scroll}の正規化母線族は
$K$ 自体に降りる。実際各一般入力平面との交点を与える写像は、
$G\to E$ を使って $K$ に降りている。二つの一般平面の共通直線は
$K$ の特異点の上の有限本の母線を避けるので、二つの写像の値はその付近で相異なる。
そこで二点を結ぶ直線を取ると、滑らかな部分の母線写像と貼り合う。

入力四空間を $V$、$K$ 上の母線部分二束を $U\subset V\otimes\mathcal O_K$、
商二束を $F$ とする。$S=\mathbb P_K(U)$ の入力射 $f$ は有限である
（正規化母線族が $S$ に有限全射し、その入力射が有限）。
核評価から $M_D=\operatorname{coker}N\hookrightarrow f_*\pi^*P$ で、
$\deg F=\deg U^*=8$、$\deg P=5$、$p_a(K)=2$ より右辺のHilbert多項式は
$(t+1)(4t+5+1-2)=4(t+1)^2$。よって同型である。
各 $[v]\in\mathbb P(V)$ での $F$ の切断の零スキームは入力繊維で、
その長さは $4-\operatorname{rank}N(v)$ に等しい。
定理\ref{thm:rank-one-skew-pair-exclusion}と命題\ref{prop:degree-six-rank-one-component-partitions}の
直線・二次曲線上の印数評価により、曲線上の印は高々六個である
（直線三個、整二次六個、二本の直線なら合わせて六個）。
従って少なくとも十四個の印が階数一の孤立点である。
それぞれは長さ三の部分スキームをただ一つ含むので
$\operatorname{Hilb}^3(K)\times\mathbb P(V)$ 上の三点繊維の孤立点を与える。

\end{proof}
\begin{proof}[曲線と四切断の変形]
まず $\operatorname{Hom}(F,P)=0$ である。非零な $F\to P$ を表示の四つの切断
$V\to H^0(F)$ と合成すれば $H^0(P)\otimes V^*$ の非零元となる。
母線上では $v$ の $F$ での値が零なので、これは $f_*\pi^*P$ の一次関係を与える。
しかし $M_D$ の表示には一次関係がない。$\deg(P\omega_K^{-1})=3$ なので
Riemann--Rochから $\omega_K\hookrightarrow P$ があり、$\operatorname{Hom}(F,\omega_K)=0$。
$U\subset V\otimes\mathcal O_K$ と双対性により $H^1(U^*\otimes F)=0$ および $H^1(F)=0$ となる。

$K$ は二重被覆として平滑化できる。実際 $\omega_K$ の二切断は基点を持たない。
任意の点 $p$ で $h^0(\mathfrak m_p^*)=1$ を示せば双対性で従う。
非正則点では $(\mathcal O:\mathfrak m_p)=(\mathfrak m_p:\mathfrak m_p)$ である。
両者が異なるなら、極大イデアルを局所環に入れるある乗法がその元を単元に移し、
極大イデアルは主となるからである。このとき大域切断は整な固有中間曲線の関数である。
正則点で非定数切断があれば一単極を持つ関数が $K\to\mathbb P^1$ の次数一の有限射を与え、
正規標的への双有理有限射なので同型となってしまう。従って標準系は基点なしで次数二を与える。
射は有限平坦で、跡分解で代数は $\mathcal O\oplus\mathcal O(-3)$、
積は二項六次式で定まる。平方因子をもつ場合も式を一般に動かせば滑らかな繊維を得る。
この表示は $K$ が局所平面曲線であることも示す。
Grassmann射は $H^1(U^*\otimes F)=0$ により同時に変形できる
（アフィン被覆上の持上げの差は接束 $U^*\otimes F$ の一次コサイクルとなる）。

滑らかな種数二曲線 $B$ 上で、次数八、$H^1(F')=0$ の二束と
その生成順序付き四切断の族は既約なパラメータ族で結べる。
具体的には比較する両二束が十分高い同一の捻りで生成されるように選ぶ。
階数二束の一般切断は曲線上で零点を持たないため、両方とも同じ十分負の線束 $A$ を
部分束に持ち、商はある線束 $D'$ である。$D'$ は同じ次数のPicard多様体上を動き、
$\operatorname{Ext}^1(D',A)$ は $\deg A-\deg D'<0$ より一定階数のベクトル束をなす。
この拡大の族で $H^1=0$ の開部分を取り、四切断はその上の切断束の四重直和の開集合である。
従って有限な三点繊維軌跡を持つ一例があれば、任意の当該四切断はその開集合の閉包に入る。

その一例には一般の異なる次数四線束 $A_1,A_2$ の直和を取る。各々は生成され切断数三。
一般四空間は直和を生成する（固定点で生成しない条件は余次元三）。
三点で消える切断の射影像は完全切断空間 $\mathbb P^5$ 内で高々二次元である。
一方の成分が零なら射影平面内にある。両成分が非零で因子 $Z$ を共有すれば
$A_i(-Z)=\mathcal O_B(p_i)$ である。$\mathcal O_B(Z)$ は $\operatorname{Pic}^3 B$ 内の
二つのAbel曲線の共通点であり、一般の $A_i$ ではこれらは異なり共通点は有限。
各共通点では $Z$ は次数三の完全鉛筆を動き、その上で消える切断は射影直線をなす。
よって一般の $\mathbb P^3$ で切ると有限であり、各非零切断の零因子に入る長さ三の因子も有限個である。
固有性から有限性は開条件である。平滑化したGrassmann射の各滑らかな繊維でも
上の既約族を使えるので、元の特異繊維は滑らかでこの有限性を持つ開集合の閉包に入る。
ここでは平滑化上の相対Homスキームを使えばよい。その閉包内で元の点を通り
一般点がこの開集合に入る既約曲線を選び、正規化して一次元の変形を取る。
\end{proof}

\begin{proof}[孤立点と三点繊維長]
この一次元族で三点の相対Hilbertスキームと入力射影空間の積を取り、
階数六の同伴束の普遍切断を零とする。中心の十四個以上の孤立点は全て一般繊維に寄与する。
実際各長さ三スキームの台の付近で曲線族は滑らかな相対平面内の一方程式で表せる。
その相対平面の三点Hilbertスキームは滑らかで相対次元六、曲線上にある条件は高々三式である。
入力射影空間の三次元と底の一次元を加え、さらに普遍切断の六式を課しても局所次元は一以上。
分離した台では長さを配分して同じ計算となる。中心で孤立しているから水平成分が必要である。
固有射を底の厳密Hensel化で考え、Stein分解で特殊繊維のこれらの孤立連結成分を
分ければ、それぞれ一般繊維の長さに一以上を与える。中心に正次元成分があってもよい。
一方滑らかな一般繊維では有限零点スキームの長さは
系\ref{cor:kernel-curve-triple-ruling-genus}の対角積分で
\[
         (8-4)((8-2)(8-3)-6\cdot2)/6=12
\]
である。これは十四以上に反する。
\end{proof}
\path{computations/verify_genus_two_kernel_excess.py} は三点の対角積分と
各階数一曲線配置で必要な孤立点数を検算する。
